\documentclass{article}
\usepackage[a4paper, total={6.5in, 10in}]{geometry}
\setlength{\parindent}{0pt}
\usepackage{tcolorbox}
\tcbset{colback=white!94!black, colframe=black!60!white, coltitle = white, boxrule=0.8pt, arc=2mm}
\newtcolorbox{boxs}[2][]{title= #2,#1}
\usepackage{datetime2}
\usepackage[british]{babel}
\usepackage{amsfonts}
\usepackage{amsmath}

\title{\textbf{Advanced Calculus HW 9}}
\author{Robert O'Connell}
\date{\today}
\begin{document}
\maketitle
\vspace{5mm}
\subsection*{Question 1}
\subsubsection*{(a)}
\begin{align*}
    \int_{1}^{2} \int_{0}^{y} xy^2 dx dy &= \int_{1}^{2} \frac{y^4}{2} dy \\
    &= \frac{1}{10}(2^5 - 1) \\
    &= 3.1
\end{align*}

\begin{align*}
    \int_{0}^{1} \int_{1}^{2} xy^2 dy dx + \int_{1}^{2} \int_{x}^{2} xy^2 dy dx &= \int_{0}^{1} x(\frac{7}{3}) dx + \int_{1}^{2} x(\frac{8}{3} - \frac{x^3}{3}) dx \\
    &= \frac{7}{6} + \frac{16}{5} - \frac{19}{15} \\
    &= 3.1
\end{align*}

\subsubsection*{(b)}
\begin{align*}
    \int_{0}^{5} \int_{5-x}^{\sqrt{25 - x^2}} y dy dx &= \frac{1}{2} \int_{0}^{5} (25 - x^2 - 25 + 10x - x^2) dx \\
    &= \frac{125}{6}
\end{align*}

\begin{align*}
    \int_{0}^{5}\int_{5-y}^{\sqrt{25-y^2}} y dx dy &= \int_{0}^{5} y(\sqrt{25 - y^2} - 5+y) dy \\
    &= \left[ -\frac{1}{3}(25-y^2)^\frac{3}{2} - \frac{5}{2}y^2 + \frac{y^3}{3}\right]_0^5 \\
    &= \frac{125}{6}
\end{align*}

\subsection*{Question 2}
Note, the plane \(x+z =1\) intersects the other plane at the line \(y = 1\), 

For the second plane, \(y \geq 0\) and \(y \leq 1\) in the region before \(y = 1\)
\begin{align*}
    \int_{0}^{1} \int_{-\sqrt{x}}^{\sqrt{x}} (1-x) dy dx &= \int_{0}^{1} 2(\sqrt{x} - x^{\frac{3}{2}}) dx \\
    &= \frac{8}{15}
\end{align*}

\subsection*{Question 3}
\begin{align*}
    \int_{0}^{1} \int_{0}^{2x} \cos(x^2) dy dx &= \int_{0}^{1} 2x \cos(x^2) dx \\
    &= \left[\sin(x^2)\right]_0^1 \\
    &= \sin(1)
\end{align*}

\subsection*{Question 4}
\(\sqrt{1-x^2 - y^2}\) is a unit hemisphere

The region we are integrating over is a unit circle in the first quadrant

Hence the volume above this region is simply an eigth of the volume of a unit sphere
\[
\frac{1}{8}\left(\frac{4}{3}\pi\right) = \frac{\pi}{6}
\]




\end{document}